%! AP Statistics — Unit 4, Lesson 4.4 — Guided Notes KEY
\documentclass[10pt]{article}
\usepackage{apstats-article}
\usepackage{apstats-key}

\begin{document}

\pageheader{Unit 4, Lesson 4.4}{Guided Notes --- KEY}
\namedateperiod

\vspace{0.12in}

\begin{objectivebox}
\small By the end of this lesson, I will be able to\dots\\[4pt]
\ansline{Define joint probability and the intersection $A \cap B$.}\\[1pt]
\ansline{Determine whether two events are mutually exclusive.}\\[1pt]
\ansline{Justify in context whether two events are mutually exclusive.}
\end{objectivebox}

\vspace{0.1in}

\begin{vocabbox}
\noindent\textbf{\textcolor{navy}{Joint probability $P(A \cap B)$:}}\\[1pt]%
\ansline{The probability that both $A$ and $B$ occur.}\\[1pt]\ansline{}\\[3pt]
\noindent\textbf{\textcolor{navy}{Intersection $A \cap B$:}}\\[1pt]%
\ansline{All outcomes that belong to both event $A$ and event $B$.}\\[1pt]\ansline{}\\[3pt]
\noindent\textbf{\textcolor{navy}{Mutually exclusive (disjoint) events:}}\\[1pt]%
\ansline{Events that cannot occur at the same time (no shared outcomes),}\\[1pt]\ansline{so $P(A \cap B) = 0$.}\\[3pt]
\noindent\textbf{\textcolor{navy}{Sample space:}}\\[1pt]%
\ansline{The set of all possible outcomes of a chance process.}\\[1pt]\ansline{}\\[3pt]
\end{vocabbox}

\vspace{0.1in}

\begin{hookbox}
\small
\textbf{Two scenarios:}
\begin{enumerate}[leftmargin=1.5em, itemsep=4pt]
  \item A randomly selected student is enrolled in both Statistics and Calculus.
  \item A randomly selected card from a standard deck is both a heart and a spade.
\end{enumerate}
\textbf{Which is impossible? Why?}\\[4pt]
\ansline{Scenario 2. A card has exactly one suit, so it cannot be both a heart and a spade.}\\[1pt]\ansline{Scenario 1 is possible: a student can take both courses.}
\end{hookbox}

\vspace{0.1in}

\begin{notesbox}{Joint Probability and Intersections}
\small
\textbf{Definition:} The \textbf{joint probability} $P(A \cap B)$ is the probability that
event $A$ \ans{occurs} and event $B$ \ans{occurs} simultaneously.

\vspace{10pt}
\textbf{Notation:} $A \cap B$ is read as ``$A$ \ans{and} $B$'' and represents all
outcomes in \ans{both $A$ and $B$}.

\vspace{10pt}
\textbf{Two Venn Diagrams:}

\vspace{6pt}
\begin{minipage}[t]{0.46\linewidth}
  \centering
  \begin{tcolorbox}[colback=white, colframe=navy, arc=2mm, width=\linewidth,
      left=2mm, right=2mm, top=2mm, bottom=2mm]
    \centering\small
    \textbf{Not Mutually Exclusive}\\[4pt]
    \begin{tikzpicture}
      \draw[navy,thick] (0,0) circle (0.7cm) node[left=0.55cm,navy]{\small$A$};
      \draw[navy,thick] (0.9,0) circle (0.7cm) node[right=0.55cm,navy]{\small$B$};
      \fill[sky,opacity=0.5] (0.45,0) ellipse (0.26cm and 0.55cm);
      \node at (0.45,0) {\tiny$A\cap B$};
    \end{tikzpicture}\\[4pt]
    $P(A \cap B)$ \ans{$> 0$}
  \end{tcolorbox}
\end{minipage}
\hfill
\begin{minipage}[t]{0.46\linewidth}
  \centering
  \begin{tcolorbox}[colback=white, colframe=navy, arc=2mm, width=\linewidth,
      left=2mm, right=2mm, top=2mm, bottom=2mm]
    \centering\small
    \textbf{Mutually Exclusive}\\[4pt]
    \begin{tikzpicture}
      \draw[navy,thick] (-0.85,0) circle (0.7cm) node[navy]{\small$A$};
      \draw[navy,thick] (0.85,0) circle (0.7cm) node[navy]{\small$B$};
    \end{tikzpicture}\\[4pt]
    $P(A \cap B) = $ \ans{$0$}
  \end{tcolorbox}
\end{minipage}
\end{notesbox}

\vspace{0.1in}

\begin{notesbox}{Mutually Exclusive Events}
\small
\textbf{Definition (VAR-4.C.2):} Two events $A$ and $B$ are \textbf{mutually exclusive}
(also called \textbf{disjoint}) if \ans{they share no outcomes}.\\[6pt]
Therefore, $P(A \cap B) = $ \ans{$0$}.

\vspace{10pt}
\textbf{How to check:}
\begin{enumerate}[leftmargin=1.5em, itemsep=4pt]
  \item List the outcomes in event $A$: \ans{(every outcome in $A$)}
  \item List the outcomes in event $B$: \ans{(every outcome in $B$)}
  \item Ask: does any single outcome appear in \ans{both $A$ and $B$}?\\
        If YES $\to$ \ans{NOT mutually exclusive} \quad If NO $\to$ \ans{mutually exclusive}
\end{enumerate}

\vspace{10pt}
\textbf{Important:} Mutually exclusive $\neq$ \ans{independent} (coming in Lesson 4.6!).
\end{notesbox}

\vspace{0.1in}

\begin{notesbox}{Worked Examples}
\small
\textbf{Example 1:} Roll a fair die. $A = \{\text{odd}\}$, $B = \{\text{even}\}$.
\begin{itemize}[itemsep=3pt]
  \item Outcomes in $A$: \ans{$\{1, 3, 5\}$}
  \item Outcomes in $B$: \ans{$\{2, 4, 6\}$}
  \item $A \cap B = $ \ans{$\varnothing$} \quad So: mutually exclusive? \ans{Yes}
  \item $P(A \cap B) = $ \ans{$0$}
\end{itemize}

\vspace{8pt}
\textbf{Example 2:} Draw a card. $A = \{\text{heart}\}$, $B = \{\text{face card}\}$.
\begin{itemize}[itemsep=3pt]
  \item Outcomes in both: \ans{Jack, Queen, King of hearts} (name them)
  \item $A \cap B = \varnothing$? \ans{No} \quad So: mutually exclusive? \ans{No}
  \item $P(A \cap B) = $ \ans{$3/52$}
\end{itemize}

\vspace{8pt}
\textbf{AP Exam tip:} Always write a complete sentence justification:\\[4pt]
``Events $A$ and $B$ \ans{are / are not} mutually exclusive because \ans{(shared outcome or none)}.''
\end{notesbox}

\vspace{0.1in}

\begin{practicebox}
\small For each pair, (1) list any shared outcomes, (2) state whether the events are
mutually exclusive, (3) justify your answer in context.

\vspace{6pt}
\begin{enumerate}[leftmargin=1.5em, itemsep=10pt]
  \item A student is randomly selected. $A = \{\text{sophomore}\}$, $B = \{\text{on the honor roll}\}$.\\
        Shared outcomes: \ans{honor-roll sophomores} \quad ME? \ans{No}\\
        Justification: \ansline{A sophomore can be on the honor roll; both can happen at once.}

  \item A month is randomly chosen. $A = \{\text{has 31 days}\}$, $B = \{\text{is in summer}\}$.\\
        Shared outcomes: \ans{July, August} \quad ME? \ans{No}\\
        Justification: \ansline{July and August have 31 days and are in summer.}

  \item An integer from 1--10 is chosen. $A = \{\text{multiple of 3}\}$, $B = \{\text{multiple of 5}\}$.\\
        Shared outcomes: \ans{None} \quad ME? \ans{Yes}\\
        Justification: \ansline{No integer 1--10 is a multiple of both 3 and 5 (15 is too big).}
\end{enumerate}
\end{practicebox}

\end{document}
